\chapter{Migração Legal para o Paraguai}\label{migracao-legal}

O Paraguai possui um dos regimes de residência mais acessíveis da América do Sul para estrangeiros. O processo é gerenciado pela \textbf{Dirección General de Migraciones} (DGM), vinculada ao Ministerio del Interior, com sede em Asunción. Este capítulo descreve o processo do ponto de vista prático, com requisitos, etapas, prazos e custos atualizados a 2024.

\medskip
\textit{Aviso legal: este capítulo é informativo e não substitui assessoria jurídica especializada. A legislação migratória pode ser alterada; verifique sempre os requisitos atualizados diretamente com a DGM ou com um advogado paraguaio habilitado.}

\section*{Tipos de Residência}
\addcontentsline{toc}{section}{Tipos de Residência}

\small
\begin{tabularx}{\linewidth}{p{2.8cm}p{2.0cm}Xp{2.0cm}}
\toprule
\textbf{Modalidade} & \textbf{Duração} & \textbf{Requisito Principal} & \textbf{Trabalhar?} \\
\midrule
Residência Temporária & 2 anos (renovável) & Propósito declarado & Sim* \\
Residência Permanente & Indefinida & 3 anos de temporária \textit{ou} econômico & Sim \\
Residência por Investimento & Direta & USD 75.000 em empresa & Sim \\
Reagrupamento Familiar & Temp. → Perm. & Cônjuge/filhos de residente & Sim \\
Aposentado & Direta & USD 1.500/mês renda & Não* \\
\bottomrule
\end{tabularx}
\vspace{-0.3em}\footnotesize *Com permissão específica / Não remunerado.\normalsize

\section*{Residência Temporária --- O Caminho Mais Comum}
\addcontentsline{toc}{section}{Residência Temporária}

A residência temporária é o primeiro passo para a maioria dos estrangeiros. Após 3 anos contínuos, é possível solicitar a residência permanente.

\subsection*{Documentos Necessários}

\begin{enumerate}
\item \textbf{Passaporte válido} com no mínimo 6 meses de validade e fotocópias de todas as páginas.
\item \textbf{Certidão de nascimento} apostilada e traduzida por tradutor juramentado para o espanhol.
\item \textbf{Certidão de antecedentes criminais} do país de origem, apostilada, com no máximo 90 dias de emissão.
\item \textbf{Certificado médico} emitido por médico habilitado no Paraguai, confirmando aptidão física e ausência de doenças infectocontagiosas (realizado presencialmente em Assunção).
\item \textbf{2 fotos 3×4} recentes (fundo branco).
\item \textbf{Comprovante de meios de subsistência}: extrato bancário, contrato de trabalho paraguaio, contrato de arrendamento rural, OU declaração de investimento.
\item \textbf{Contrato de locação ou escritura} de imóvel no Paraguai (domicílio comprovado).
\item \textbf{Formulário de solicitação} preenchido e assinado perante a DGM.
\end{enumerate}

\textit{Para casados:} acrescentar certidão de casamento apostilada e traduzida. Para filhos menores acompanhantes: certidão de nascimento de cada filho.

\subsection*{Etapas do Processo}

\begin{enumerate}
\item \textbf{Apostila e tradução} (país de origem): apostile a certidão de nascimento e de antecedentes criminais segundo a Convenção de Haia. Contrate tradução juramentada para espanhol.
\item \textbf{Entrada no Paraguai}: pode entrar como turista (90 dias prorrogáveis), enquanto o processo tramita.
\item \textbf{Protocolo na DGM}: compareça pessoalmente à DGM em Assunção (Av. República Argentina 537, ou escritórios regionais) com todos os documentos. Receba número de protocolo.
\item \textbf{Exame médico}: realize o exame no Centro Médico da DGM ou clínica credenciada.
\item \textbf{Análise}: a DGM analisa o processo em 30--90 dias úteis (podendo variar conforme volume de solicitações).
\item \textbf{Aprovação e cédula de identidade}: com a residência aprovada, solicite a \textit{cédula de identidad} paraguaia na Policía Nacional.
\end{enumerate}

\subsection*{Custos Estimados (2024)}

\begin{tabularx}{\linewidth}{lX}
\toprule
\textbf{Item} & \textbf{Custo Estimado (USD)} \\
\midrule
Taxa DGM --- processamento & 130--180 \\
Exame médico DGM & 20--40 \\
Cédula de identidade paraguaia & 10--15 \\
Apostila (país de origem, por documento) & 20--80 \\
Tradução juramentada (por documento) & 30--80 \\
Assessoria de advogado/despachante & 300--800 \\
\midrule
\textbf{Total estimado} & \textbf{500--1.200} \\
\bottomrule
\end{tabularx}

\section*{Residência Permanente}
\addcontentsline{toc}{section}{Residência Permanente}

Após 3 anos de residência temporária contínua com entradas e saídas regulares documentadas, o estrangeiro pode solicitar a residência permanente. O processo é similar ao da temporária: atualização dos documentos (certidão de antecedentes atualizada, exame médico, comprovante de domicílio) e comparecimento presencial à DGM.

\textbf{Via investimento direto (USD 75.000):} possível solicitar a residência permanente diretamente, sem passar pela temporária, mediante comprovação de investimento em empresa constituída no Paraguai (capital integralizado e documentado). Exige apresentação de Contrato Social registrado, comprovante bancário do aporte e certidão da SET (\textit{Subsecretaría de Estado de Tributación}).

\section*{Naturalização}
\addcontentsline{toc}{section}{Naturalização}

O Paraguai admite naturalização após \textbf{3 anos de residência permanente} (totalizando, na prática, 6 anos desde a entrada). O processo tramita no Poder Legislativo (Senado da República) mediante solicitação ao Ministerio de Relaciones Exteriores. Requisitos adicionais:

\begin{itemize}
\item Conhecimento básico de espanhol ou guarani demonstrado em entrevista.
\item Renúncia à cidadania anterior (o Paraguai \textit{não} admite dupla cidadania como regra geral, mas alguns países aceitam a perda formal e a reconhecem posteriormente --- verifique a legislação do seu país de origem).
\item Registro no \textit{Registro Cívico Permanente} (TSJE) para exercício do voto.
\end{itemize}

\section*{Contatos e Endereços Oficiais}
\addcontentsline{toc}{section}{Contatos Oficiais}

\begin{tabularx}{\linewidth}{lX}
\toprule
\textbf{Órgão} & \textbf{Endereço / Contato} \\
\midrule
Dirección General de Migraciones (DGM) & Av. República Argentina 537 e O'Leary, Asunción. Tel: +595 21 446-020. \textit{migraciones.gov.py} \\
Policía Nacional --- Cédulas & Carlos Antonio López 3547, Asunción. Horário: seg--sex 7h--15h. \\
SET (tributação) & Chile 1089 c/ Palma, Asunción. Tel: +595 21 417-7000. \textit{set.gov.py} \\
DGEEC (estatísticas, Censo 2022) & Presidente Franco 1661, Asunción. \textit{dgeec.gov.py} \\
INDERT (terras rurais) & Gral. Díaz 521, Asunción. Tel: +595 21 610-1400. \textit{indert.gov.py} \\
SENAVE (agropecuário, importação) & Ruta Mcal. Estigarribia km 10,5, San Lorenzo. \textit{senave.gov.py} \\
\bottomrule
\end{tabularx}

\section*{Abertura de Empresa no Paraguai}
\addcontentsline{toc}{section}{Abertura de Empresa}

Estrangeiros com residência temporária ou permanente podem abrir empresa no Paraguai. O tipo societário mais comum para estrangeiros é a \textbf{SRL} (\textit{Sociedad de Responsabilidad Limitada}) ou a \textbf{SA} (\textit{Sociedad Anónima}).

\begin{tabularx}{\linewidth}{lX}
\toprule
\textbf{Parâmetro} & \textbf{Valor} \\
\midrule
Capital mínimo SRL & PYG 1 (sem mínimo legal prático) \\
Prazo médio de constituição & 5--15 dias úteis \\
Custo de constituição (honorários + taxas) & USD 500--1.500 \\
Imposto de Renda Empresarial (IRACIS) & 10\% sobre lucro \\
IVA (Impuesto al Valor Agregado) & 10\% (geral) / 5\% (produtos básicos) \\
Distribuição de dividendos ao exterior & 15\% de retenção na fonte \\
\bottomrule
\end{tabularx}

O Paraguai adota o \textbf{princípio da territorialidade}: apenas rendas geradas \textit{dentro} do Paraguai são tributadas. Rendimentos de fontes estrangeiras não são tributados localmente.

\section*{Contas Bancárias}
\addcontentsline{toc}{section}{Contas Bancárias}

Estrangeiros com residência válida podem abrir conta em qualquer banco paraguaio. Os documentos típicos exigidos são: cédula de identidade paraguaia (ou passaporte + comprovante de residência), comprovante de renda ou referência bancária do exterior, e formulário interno do banco (KYC/AML).

Bancos com maior presença nacional: \textbf{Banco Continental}, \textbf{BBVA Paraguay}, \textbf{Itaú Paraguay}, \textbf{Banco Nacional de Fomento (BNF)} (público, crédito rural). Cooperativas de crédito (\textit{cooperativas de ahorro y crédito}) são alternativa acessível no interior, com taxas competitivas e menor burocracia.

\section*{Dicas Práticas}
\addcontentsline{toc}{section}{Dicas Práticas}

\begin{itemize}
\item \textbf{Contrate um despachante ou advogado local}: o processo é gerenciável sem assessoria, mas um profissional local economiza tempo e evita erros de documentação. Custo típico: USD 300--500 para residência simples.
\item \textbf{Apostile tudo antes de viajar}: apostila de documentos emitidos fora do Paraguai deve ser feita no país de origem. No Brasil, cartórios conveniados ao CNJ emitem apostila. Em outros países, verifique o órgão equivalente.
\item \textbf{Não fique em situação irregular}: ultrapassar 90 dias como turista sem protocolar a residência gera multa e pode complicar o processo. O visto de turista pode ser prorrogado na DGM por mais 90 dias.
\item \textbf{Aprenda espanhol básico}: embora muitos paraguaios no interior falem guarani como língua principal, o espanhol é o idioma de todos os trâmites oficiais. Não há tradução oficial para inglês ou português nos órgãos públicos.
\item \textbf{Guarani (idioma)}: o guarani é cooficial. Em zonas rurais e comunidades indígenas, conhecer guarani básico é diferencial significativo de integração.
\item \textbf{RUC (Registro Único de Contribuyentes)}: necessário para qualquer atividade econômica formal. Obtido na SET com a cédula paraguaia.
\end{itemize}

\clearpage
